\documentclass{article}
\usepackage[utf8]{inputenc}
\usepackage{amsmath}
\usepackage{geometry}
\geometry{margin=2cm}
\begin{document}

\section*{Questão 92 -- ENEM 2024 -- Ciências da Natureza}

\subsection*{Enunciado}
O magnésio metálico utilizado em ligas leves \`e produzido em um processo que envolve várias etapas e utiliza \`agua do mar como matéria-prima. A primeira etapa desse processo consiste na reação entre o \i{ion} $\mathrm{Mg^{2+}}$ e hidróxido de cálcio, $\mathrm{Ca(OH)_2}$, obtendo uma mistura que contém hidróxido de magnésio, pouco solúvel, e \i{ons} $\mathrm{Ca^{2+}}$, de acordo com a equação química:

\begin{equation*}
\mathrm{Mg^{2+} (aq) + Ca(OH)_2 (aq) \rightarrow Mg(OH)_2 (s) + Ca^{2+} (aq)}
\end{equation*}

O método adequado para separar o $\mathrm{Mg(OH)_2}$ dessa mistura \`e a:

\begin{enumerate}
\item Filtração
\item Catação
\item Destilação
\item Dissolução
\item Evaporação
\end{enumerate}

\section*{Interpretação e Estratégia}
\noindent A questão solicita a identificação do método apropriado para separar o sólido pouco solúvel $\mathrm{Mg(OH)_2}$ da mistura aquosa que contém \i{ons} $\mathrm{Ca^{2+}}$.

\paragraph{Termos do enunciado:}
\begin{itemize}
\item \textbf{Separar}: isolar ou dividir os componentes da mistura.
\item \textbf{Método}: procedimento ou técnica utilizada para realizar a separação.
\end{itemize}

\paragraph{O que a questão pede:}
Indicar o método adequado para separar o hidróxido de magnésio sólido suspenso em solução aquosa contendo \i{ons} cálcio, após a reação química dada.

\paragraph{Estratégia:}
Compreender os estados físicos dos reagentes e produtos. O $\mathrm{Mg(OH)_2}$ está em sólido pouco solúvel e o $\mathrm{Ca^{2+}}$ permanece em solução. Selecionar o método físico eficaz para separar um sólido em suspensão de uma solução líquida, descartando técnicas inadeqúadas.

\section*{Conteúdo e Mini Revisão}
\begin{tabular}{|p{5cm}|p{6cm}|p{5cm}|}
\hline
\textbf{Conceito} & \textbf{Ideia-chave} & \textbf{Aplicacão} \\
\hline
Sólido pouco solúvel em suspensão & Pode ser separado da solução por métodos físicos que isolam sólidos de líquidos & Filtração \`e método prático e eficaz para separar suspensões de sólidos \\
\hline
Estados físicos dos componentes & Identificar sólido, líquido e \i{ons} dissolvidos para escolher técnica correta & $\mathrm{Mg(OH)_2}$ sólido e $\mathrm{Ca^{2+}}$ em solução permitem filtração \\
\hline
Métodos físicos de separação & Escolher método que não altere composição química ao separar fases físicas distintas & Filtração separa sólido sem alterar substâncias \\
\hline
\end{tabular}

\paragraph{Teoria:}
\begin{itemize}
\item Em misturas onde um sólido pouco solúvel está suspenso em solução aquosa, o método mais comum para separar o sólido \`e a filtração, que utiliza um meio filtrante para reter o sólido enquanto deixa passar o líquido.
\item A composição química e os estados físicos guiam a escolha do método. Aqui, $\mathrm{Mg(OH)_2}$ está sólido, enquanto $\mathrm{Ca^{2+}}$ permanece em solução.
\item Outras técnicas, como catação, são inadequadas para sólidos microscópicos dispersos; destilação separa líquidos; e evaporação \`e menos prática para esta situação.
\end{itemize}

\section*{Resolução e Pulo do Gato}
\begin{enumerate}
\item Observe a reação: $\mathrm{Mg^{2+} (aq) + Ca(OH)_2 (aq) \rightarrow Mg(OH)_2 (s) + Ca^{2+} (aq)}$. O $\mathrm{Mg(OH)_2}$ surge como sólido pouco solúvel em suspensão, e o $\mathrm{Ca^{2+}}$ permanece em solução.
\item Para separar o $\mathrm{Mg(OH)_2}$, utiliza-se a filtração, que retém o sólido e permite passar o líquido com os \i{ons}.
\item A alternativa correta \`e a \textbf{A}, filtração.
\item Outros métodos: catação \`e manual e para sólidos visíveis; destilação \`e para misturas líquidas; dissolução mistura, não separa; evaporação \`e mais lenta e energética.
\end{enumerate}

\paragraph{Pulo do gato:}
\textit{Identifique o estado físico dos componentes e lembre-se: para separar sólido pouco solúvel em líquido, o método adequado \`e a filtração, simples e eficiente.}

\section*{Análise dos Distratores}
\begin{itemize}
\item \textbf{B}: Catação confundida com filtração. Catação não separa sólido microscópico em líquido.
\item \textbf{C}: Destilação serve para separar líquidos, não sólidos suspensos.
\item \textbf{D}: Dissolução mistura componentes, não separa.
\item \textbf{E}: Evaporação \`e possível, mas menos prática que filtração e mais demorada.
\end{itemize}

\section*{Código da Questão}
\texttt{CIENCIAS_DA_NATUREZA_QUIMICA_001}

\end{document}